\section{What the chaos does to inference and steering}
\label{sec:consequences}
The loop studied here is the expected-embedding feedback that continuous reasoning methods such as Soft Thinking use
\cite{zhang2025softthinking}, and its chaos has consequences that can be measured without any reference to dynamical
systems. We report three: whether a run can be reproduced, how the depth at which a steering vector enters sets the
chaos it causes, and what happens when two steering vectors are combined. A fourth, that a repeating text can hide a
chaotic state from anyone who reads only the output, is shown in Section~\ref{sec:hostility}.

\paragraph{The regime decides whether a run can be reproduced.} Language-model inference is not bitwise
reproducible: GPU kernels choose different reduction orders for different matrix shapes, so a request's result depends
on the batch it shares and on the hardware \cite{he2025nondeterminism,yuan2025nondeterminism}. We ran the 448 labelled
cells of the horizon check (192 chaotic, 192 periodic and 64 other cells of two float64 maps; Section~\ref{sec:methods})
twice in lockstep from the same start. The two runs differ only in the batch: all cells in one batch of $n$ rows, or
the same rows followed by the same cells in reverse order ($2n$ rows), which changes the matrix shapes the kernels see
and nothing else (Figure~\ref{fig:repro},
Table~\ref{tab:repro}). In float64 the two runs first differ by about $6\times10^{-15}$, the rounding level. At 190 of
the 192 chaotic cells this difference grows to order one, after a median of 102 steps, and 90\,\% of the chaotic cells
end up writing different tokens. At 182 of the 192 periodic cells the two runs stay within $10^{-2}$ for all 1500
steps; their median largest distance is $2\times10^{-12}$, and 95\,\% of all periodic cells end the run within
$3\times10^{-12}$. The Lyapunov exponent predicts when a chaotic run
diverges: the step at which the distance passes $10^{-2}$ follows $\ln(10^{-2}/d_1)/\lambda$, where $d_1$ is the first
nonzero distance, with a median ratio of $0.91$ and a log-log correlation of $0.73$ over 190 cells. The two chaotic
cells that never diverge have exponents of $0.013$ and $0.0013$, for which the predicted step lies beyond the run.
The same holds without steering in Qwen2.5-7B, whose embeddings are untied. Over 24 chaotic and 24 periodic
temperatures of its unsteered loop, all 24 chaotic float64 runs diverge, after a median of 56 steps (the exponent again
predicts the step, median ratio $0.95$), and 96\,\% of them end up writing different tokens. None of the 24 periodic runs
moves past $10^{-2}$, and 95\,\% end within $1.3\times10^{-13}$ of each other. In lower precision the rounding is
larger. In float32 every chaotic cell of both models diverges, after a median of 20 steps (0.5B) or 16 (7B). In
bfloat16, a common serving precision, the two layouts pass a distance of $10^{-2}$ after a median of two steps, at every
cell, periodic ones included. At that precision the loop is not reproducible across batch layouts in any regime we
tested.

The same test at the level of answers, with a design frozen before the run, shows what this means for a task. We
asked Qwen2.5-0.5B 96 held-out two-digit additions through a short soft-thinking pipeline: 32 expected-embedding steps
appended to a growing context, then a fixed answer cue and greedy extraction of the number. Each problem ran alone and
in a batch with three fixed distractor problems of the same length, so nothing is padded and only the matrix shapes
differ. In bfloat16 the extracted answer changed with the batch layout in 56\,\% of the problems at $T=1$ and in
8\,\% at $T=10^{0.2}$. Ordinary greedy decoding, run with the same budget, produced identical token sequences in both
layouts for every problem. The excess disagreement of soft over greedy decoding, averaged over the two temperatures,
is 32.3 percentage points (95\,\% interval $[26.6,38.0]$; sign-flip $p=10^{-4}$, the smallest attainable). In float64 no
answer changed within the 32 steps. There, a perturbation of the first soft token grows by only $0.03$ nats per step
at $T=1$ and shrinks at higher temperature, and even a relative kick of $10^{-4}$ changes no answer. The bfloat16
discrepancies are therefore not the float64 map's chaos. They arise from bfloat16's rounding differences accumulating
through continuous feedback, which a discrete argmax would absorb. (The greedy baseline's accuracy is not
informative: forced through 32 tokens, the model writes past its answer, and the cue then extracts another number.)
Two horizons thus matter. Within a short budget, continuous feedback in low precision makes answers depend on the
batch. Over longer runs, a chaotic regime amplifies even float64 rounding into different text, on the schedule its
Lyapunov exponent sets.
\begin{figure}[tbp]
\centering
\includegraphics[width=\linewidth]{fig-repro.png}
\caption{Only the batch layout differs. The 448 labelled cells of Qwen2.5-0.5B-Instruct (philosophy direction, layer
12) run twice from the same start, alone in a batch or followed by the same cells reversed. (a) float64: median and
interquartile range of the relative distance between the two runs, by the cells' float64 labels. (b) Medians for
float32 and bfloat16 batch layouts. (c) float64 chaotic cells: the step at which the distance passes $10^{-2}$ against
$\ln(10^{-2}/d_1)/\lambda$, with $d_1$ the first nonzero distance and $\lambda$ the map exponent; dotted, equality.}
\label{fig:repro}
\end{figure}
\begin{table}[tbp]
\centering\small
\begin{tabular}{llrrr}
\toprule
comparison & cells & reach $10^{-2}$ & median step & end with different tokens \\
\midrule
batch layout, float64 & chaotic & 190 / 192 & 102 & 90.1\,\% \\
 & periodic & 10 / 192 & & 3.6\,\% \\
batch layout, float32 & chaotic & 192 / 192 & 20.5 & 92.2\,\% \\
 & periodic & 72 / 192 & & 22.9\,\% \\
float32 against float64 & chaotic & 192 / 192 & 20 & 92.7\,\% \\
 & periodic & 75 / 192 & & 19.3\,\% \\
batch layout, bfloat16 & chaotic & 192 / 192 & 2 & 99.5\,\% \\
 & periodic & 192 / 192 & 2 & 95.3\,\% \\
\midrule
7B unsteered, float64 & chaotic & 24 / 24 & 56 & 95.8\,\% \\
 & periodic & 0 / 24 & & 0\,\% \\
7B unsteered, float32 & chaotic & 24 / 24 & 16 & 95.8\,\% \\
 & periodic & 4 / 24 & & 8.3\,\% \\
7B unsteered, bfloat16 & chaotic & 24 / 24 & 2 & 75.0\,\% \\
 & periodic & 24 / 24 & 2 & 50.0\,\% \\
\bottomrule
\end{tabular}
\caption{Divergence of two otherwise identical runs, by the cells' float64 labels (1500 steps; ``end with different
tokens'': the most probable tokens differ at some step of the last 500). Top: Qwen2.5-0.5B, philosophy direction at
layer 12. Bottom: Qwen2.5-7B unsteered, 24 chaotic and 24 periodic temperatures of its long sweep. In float32 the periodic cells that diverge may
change phase on the same cycle or switch between coexisting attractors; the run does not distinguish the two.}
\label{tab:repro}
\end{table}

\paragraph{Where a steering vector enters decides how much chaos it causes.} Steering vectors are a common tool for
controlling and probing models \cite{zou2023representation,rimsky2024steering}, usually added at one layer chosen by
hand. In Qwen2.5-0.5B the philosophy direction, built with the same recipe at blocks 4, 8, 12, 16 and 20 and scaled to
each block's residual norm, makes 5.9, 8.1, 6.8, 5.0 and 2.3\,\% of the same $128\times128$ plane chaotic
(Figure~\ref{fig:layersweep}). Early steering carries the chaos down to $\log_{10}T\approx-0.9$, close to greedy
decoding, while steering at block 20 confines it to a thin band at $\log_{10}T=0.1$--$0.25$. Near zero steering every
block gives 1--2\,\%. Llama-3.3-70B behaves the same way. Steered along its hostility direction, it is chaotic in 8.0\,\%
of a $24\times24$ map at block 20 of 80, down to $\log_{10}T=-0.14$, against 4.3\,\% at block 35, where the chaos starts
at $\log_{10}T=0.01$, 4.3\,\% at block 50 (32 by 32 cells), and 0.7\,\% at block 65, all above
$\log_{10}T=0.19$. These 70B labels are those of the map
protocol, one twin offset, and so are candidates. A logit bias, the latest possible point of entry, gives chaos mostly
above $\log_{10}T=0.1$ (Section~\ref{sec:theory}). The gain split of Section~\ref{sec:theory}, repeated along
steering lines at all five blocks of the 0.5B model, shows where the depth enters. At low temperature
($\log_{10}T\le0$) the chaotic cells of blocks 4, 8 and 12 have expanding readout factors (medians $+0.74$, $+0.46$ and
$+0.46$ nats per step over 54, 56 and 21 cells) and contracting network factors ($-0.47$ to $-0.21$). At block 16
the readout factor of such cells is only $+0.15$, and block 20 gives just 6 of them. Early injection creates token
competition at low temperature more readily; late injection mostly does not. Its chaos lives at higher temperature,
where the network factor carries it (block 20: network $+1.09$, readout $-0.78$; Figure~\ref{fig:mechanism}e).
\begin{figure}[tbp]
\centering
\includegraphics[width=\linewidth]{fig-layer-sweep.png}
\caption{The same steering direction (philosophy, Qwen2.5-0.5B-Instruct) entering at five blocks of 24: largest
Lyapunov exponent on $128\times128$ float64 cells over the steering strength and $\log_{10}T\in[-1,0.3]$. The share of
chaotic cells is given above each panel.}
\label{fig:layersweep}
\end{figure}

\paragraph{Two steering vectors interact.} Steering at two blocks at once, gratitude at block 4 and hostility at block
16 of Qwen2.5-0.5B, makes 21.7\,\% of a $96\times96$ plane chaotic at $\log_{10}T=-0.1$. Along the two axes, where only one
push acts, the shares are 7.3 and 8.3\,\%. Large chaotic regions lie off the axes, where neither push alone is chaotic
(Figure~\ref{fig:twoblock}): in 17.1\,\% of that plane (1574 cells, 79\,\% of its chaotic off-axis cells) the combined
push is chaotic although each push alone, at the same strength, is not. Over the four planes this share is 9.5 to
17.1\,\%. With the blocks exchanged (hostility at 4, gratitude at 16) the plane is 23.4\,\% chaotic,
and at $\log_{10}T=0.1$ both pairs give 33 and 42\,\%. Each plane keeps a periodic hole around zero steering. Combining
steering vectors is therefore not additive in its effect on the loop's stability: a composition of two pushes can be
chaotic where each push alone is periodic.
\begin{figure}[tbp]
\centering
\includegraphics[width=0.85\linewidth]{fig-two-block.png}
\caption{Two steering directions at two blocks of Qwen2.5-0.5B-Instruct: largest Lyapunov exponent on $96\times96$
float64 cells over the two strengths, at $\log_{10}T=-0.1$ (left) and $0.1$ (right). Top: gratitude at block 4,
hostility at block 16; bottom: the blocks exchanged. Dotted lines mark the axes, where only one push acts.}
\label{fig:twoblock}
\end{figure}

\paragraph{Does a chaotic loop predict worse answers in ordinary decoding? Not demonstrably.} We looked for a
behavioural signature of the regimes outside the loop, in two exploratory designs, and found none that survives
controls. First, ordinary sampled answers of Qwen2.5-0.5B under steering at the strengths of chaotic and periodic map
cells (30 matched pairs; arithmetic, open questions, short programs) showed no measurable accuracy contrast. At those
strengths, steering every position turns every reply into philosophy-flavoured text, so accuracy sits at the floor in
both zones. Second, we labelled each of 200 prompts by the regime of its own unsteered loop in Qwen2.5-7B at three
temperatures and compared the prompts' ordinary sampled accuracy. Prompts whose loop is chaotic were answered far less
accurately (by 23 to 36 points among prompts the model answers correctly when greedy), but mostly because chaotic
labels fall on the harder task and on flatter answer distributions. With task and answer entropy as covariates, a
residual of $-4$ to $-5$ points remains, with the same sign at each temperature and no significance alone. Answer
entropy is measured during sampling, so this adjustment may over-correct. A regime label measured on the loop may yet
add predictive value beyond simpler difficulty measures, but establishing it needs a frozen test on fresh, single-task
prompts with covariates fixed before sampling; we have not run one.

